% =====================================================================
\section{Experiments}\label{sec:exp}
% =====================================================================

We evaluate Walk Attention on a controlled family of graphs in which
over-squashing is severe and \emph{tunable}, so that the contribution of each
architectural component can be isolated. All code, configurations and the
experiment logs are released with the paper.

% ---------------------------------------------------------------------
\subsection{A tunable over-squashing benchmark}
% ---------------------------------------------------------------------

\paragraph{Bottleneck-chain graphs.}
We build a family parameterised by a \emph{depth} $d$, a number of sources $K$
and a width $M$. The graph has $K$ source vertices, then $d$ consecutive
\emph{bottleneck layers} of $M$ interchangeable vertices each, and finally one
target vertex. Every source connects to every vertex of the first layer, every
vertex of layer $r$ connects to every vertex of layer $r+1$, and every vertex of
the last layer connects to the target. Because the bottleneck vertices are
interchangeable, the number of length-$(d{+}1)$ walks from a source to the target
is $M^{d}$, so the walk multiplicity into the target grows as
$n_{d+1}(\cdot,\text{target})\sim K\,M^{d}$
(Proposition~\ref{prop:count}); a message-passing GNN must compress all of them
into the target's fixed-width vector. The walk-reachability support, by contrast,
remains sparse: at the deepest range it consists of only the $K$ source--target
pairs ($5/196\approx2.6\%$ of all node pairs at depth $2$ and
$5/324\approx1.5\%$ at depth $3$ in our configuration).

\paragraph{Task.}
Each source carries a one-hot \emph{key} and a scalar \emph{payload}; the target
carries a query key. The model must output, at the target node, the key of the
source whose key matches the query---a retrieval task whose answer lies $d{+}1$
hops away, beyond the bottleneck. Accuracy is measured at the target node;
chance level is $1/K$.

\paragraph{Models.} We compare:
\begin{itemize}
  \item \textbf{GAT} \cite{velickovic2017graph}: one-hop graph attention, $d{+}1$
  layers (learned weights, one-hop support);
  \item \textbf{Walk operator} (\texttt{walkraw}): the fixed-weight multi-hop
  operator $W^{(k)}$ of Definition~\ref{def:walkop} (multi-hop support, fixed
  multiplicity weights);
  \item \textbf{Walk Attention} (ours): Algorithm~\ref{alg:wa} (multi-hop
  support, learned weights).
\end{itemize}
All models use hidden width $16$, four heads where applicable, and $d{+}1$
layers/ranges. We train with AdamW, a ten-epoch linear learning-rate warmup
followed by cosine decay, and gradient clipping; Walk Attention additionally uses
layer normalisation after each layer. We report mean and standard deviation of
target accuracy on a held-out validation split (the synthetic task has no
separate test set) over five random seeds, at problem depths $d\in\{2,3\}$ with
$K=5$, $M=4$.

% ---------------------------------------------------------------------
\subsection{Main result}
% ---------------------------------------------------------------------

Table~\ref{tab:main} reports the results, which land exactly where
Propositions~\ref{prop:collapse}--\ref{prop:wa} put them
(Remark~\ref{rem:mechanisms}). One-hop GAT degrades sharply with depth,
approaching chance: a local operator cannot propagate the signal across the
bottleneck under over-squashing. GCN and GIN, whose first aggregation is
linear, are pinned at chance by Proposition~\ref{prop:blind}---and sit there
empirically at every depth. The fixed-weight walk operator reaches the
target---its multi-hop support spans the bottleneck in a single layer---but,
weighting all reachable sources by multiplicity alone, it cannot select the
matching source and plateaus at $0.50$--$0.62$. \textbf{Walk Attention solves the
task exactly}, attaining $1.000\pm0.000$ accuracy at both depths over all five
seeds: the same multi-hop support, now with learned query--key weights, lets the
target attend selectively to the relevant source.

\begin{table}[t]
\centering
\caption{Target accuracy (mean $\pm$ std over 5 seeds) on bottleneck-chain
retrieval, $K{=}5$, $M{=}4$. Chance is $1/K=0.20$. Depth $d$ is the number of
bottleneck layers; the answer lies $d{+}1$ hops away. GCN and GIN sit at
chance, as Proposition~\ref{prop:blind} mandates (the residual
$+0.01$--$0.02$ over $1/K$ is best-epoch selection on a finite validation
split).}
\label{tab:main}
\begin{tabular}{lcc}
\toprule
\textbf{Model} & \textbf{depth $d=2$} & \textbf{depth $d=3$} \\
\midrule
GCN (one-hop, linear aggregation)  & $0.21\pm0.01$ & $0.22\pm0.01$ \\
GIN (one-hop, linear aggregation)  & $0.21\pm0.01$ & $0.21\pm0.01$ \\
GAT (one-hop attention)            & $0.45\pm0.22$ & $0.29\pm0.17$ \\
Walk operator (fixed weights)      & $0.62\pm0.12$ & $0.50\pm0.12$ \\
\textbf{Walk Attention (ours)}     & $\mathbf{1.000\pm0.000}$ & $\mathbf{1.000\pm0.000}$ \\
\bottomrule
\end{tabular}
\end{table}

\paragraph{On stability.}
Without layer normalisation and warmup, Walk Attention has a high ceiling but is
high-variance: across seeds its accuracy ranged from $0.44$ to $1.00$. The three
standard stabilisers (layer normalisation, learning-rate warmup, gradient
clipping) remove this variance entirely, yielding the $\pm0.000$ in
Table~\ref{tab:main}. The instability was therefore an optimisation artefact, not
a property of the operator.

% ---------------------------------------------------------------------
\subsection{RingTransfer}
% ---------------------------------------------------------------------

The bottleneck family is our own construction. To confirm the effect on an
\emph{independent} benchmark we use \textsc{RingTransfer}
\cite{digiovanni2023oversquashing}: a cycle of $n$ nodes in which a source must
transmit its one-hot label to the diametrically opposite target. The
parallel-path structure is explicit: the source reaches the target along the
\emph{two arcs} of the ring, two paths of equal length $n/2$. Over-squashing is
severe and tunable: every node has two neighbours, so the number of
length-$(n/2)$ walks entering the target---the messages a one-hop network must
compress after $n/2$ layers---grows as $2^{\,n/2}$ with the ring size, while
only the two arc walks carry the signal. The task requires $n/2$ hops, so a
one-hop model requires $n/2$ layers. We use five classes (chance $0.2$), hidden
width $32$, the same optimiser recipe as above (AdamW, linear warmup, cosine
decay, gradient clipping), and $n/2$ layers/ranges for all models.

Table~\ref{tab:ring} reports the results. On small rings ($n\le 10$, distance
$\le 5$) every model succeeds, as the squashing is not yet severe. From $n=14$
(distance $7$) onward the one-hop baselines degrade and become highly
seed-dependent---at $n=18$ GAT falls to $0.29\pm0.19$ and GCN to
$0.53\pm0.33$---while \textbf{Walk Attention remains at $1.000\pm0.000$
throughout}: its multi-hop support aggregates the source directly along either
arc, and the learned weights identify the labelled source. The gap widens with
ring size, i.e.\ with the severity of over-squashing, as predicted.

\begin{table}[t]
\centering
\caption{\textsc{RingTransfer} accuracy vs.\ ring size $n$ (distance $n/2$),
mean $\pm$ std over 5 seeds; chance $=0.20$. One-hop GCN/GAT collapse as the
ring---and thus the over-squashing---grows, while Walk Attention stays exact.}
\label{tab:ring}
\begin{tabular}{lcccc}
\toprule
\textbf{Model} & $n{=}6$ & $n{=}10$ & $n{=}14$ & $n{=}18$ \\
\midrule
GCN                    & $1.00\pm0.00$ & $1.00\pm0.00$ & $0.81\pm0.13$ & $0.53\pm0.33$ \\
GAT                    & $1.00\pm0.00$ & $1.00\pm0.00$ & $0.69\pm0.29$ & $0.29\pm0.19$ \\
\textbf{Walk Attention (ours)} & $\mathbf{1.00\pm0.00}$ & $\mathbf{1.00\pm0.00}$ & $\mathbf{1.00\pm0.00}$ & $\mathbf{1.00\pm0.00}$ \\
\bottomrule
\end{tabular}
\end{table}

% ---------------------------------------------------------------------
\subsection{Collapsing walks hurts}
% ---------------------------------------------------------------------

The path-algebra quotient $\kQ/I$ of Section~\ref{sec:algebra} suggests an
alternative remedy: replace the walk multiplicities $n_k(i,j)$ by the smaller
quotient dimensions $\dim(e_j(\kQ/I)_k e_i)$, de-duplicating equivalent walks
before aggregation. We implemented this (the fixed-weight operator built from the
quotient dimensions) and compared it to the unquotiented walk operator on the
same task, in a controlled ablation that changes \emph{only} the operator
(quotient vs.\ raw) and holds everything else---data, architecture, and the
optimiser recipe of Table~\ref{tab:main}---fixed. As Table~\ref{tab:quotient}
shows, \textbf{the quotient operator falls clearly below the raw one at both
depths}. Collapsing parallel walks discards the source multiplicity that the
retrieval task needs: in this regime redundant paths are \emph{signal}, not
noise, and removing them removes information.

\begin{table}[t]
\centering
\caption{De-duplicating walks via the quotient $\kQ/I$ \emph{hurts}: target
accuracy (mean $\pm$ std over 5 seeds), bottleneck retrieval, under the same
protocol as Table~\ref{tab:main} (the walk-operator row coincides with it). The
quotient operator is clearly below the unquotiented one at both depths.}
\label{tab:quotient}
\begin{tabular}{lcc}
\toprule
\textbf{Operator} & \textbf{depth $d=2$} & \textbf{depth $d=3$} \\
\midrule
Quotient $\kQ/I$ (collapsed walks) & $0.39\pm0.04$ & $0.30\pm0.09$ \\
Walk operator (raw walks)          & $\mathbf{0.62\pm0.12}$ & $\mathbf{0.50\pm0.12}$ \\
\bottomrule
\end{tabular}
\end{table}

This negative result indicates that the contribution of the path algebra here is
\emph{not} compression. Its appropriate role is to supply the multi-hop attention
\emph{support}---sparsely and exactly---over which the network learns weights,
which is precisely Walk Attention. Whether a regime exists in which the
quotient's compression is itself beneficial (e.g.\ when parallel paths are
genuinely redundant noise, or when distinct path classes are routed to distinct
heads) remains an open question.

% ---------------------------------------------------------------------
\subsection{LRGB Peptides-func}
% ---------------------------------------------------------------------

The synthetic benchmark is deliberately a worst case: an extreme bottleneck where
multi-hop reach is decisive. To assess Walk Attention on real data we also
evaluate on \textbf{Peptides-func} from the Long Range Graph Benchmark
\cite{dwivedi2022lrgb} ($10{,}873/2{,}331/2{,}331$ train/val/test molecular
graphs, median $138$ nodes; multi-label graph classification, metric Average
Precision). We use the same atom embedding, four layers, hidden width $80$, mean
pooling and AdamW for every model; Walk Attention uses ranges $k=1,\dots,4$ with
the walk-reachability support of each molecule. On these graphs the support stays
sparse ($1$--$8\%$ of $n^2$ at range $3$), so the method remains cheap.

Table~\ref{tab:lrgb} reports test AP. \textbf{Walk Attention is comparable to
the one-hop baselines but does not surpass them here}: it exceeds GCN and is
statistically indistinguishable from GAT ($0.526\pm0.007$ vs.\ $0.526\pm0.004$
over four seeds). Two caveats frame this comparison. First, all models share a
deliberately small budget (four layers, width $80$, no positional encodings, a
short schedule), so our absolute numbers are below published results for this
benchmark---the LRGB paper reports GCN at $0.593$ under its tuned
$500$k-parameter protocol \cite{dwivedi2022lrgb}, and subsequent re-evaluations
are higher still \cite{tonshoff2023where}; the comparison is therefore
\emph{internal}, isolating the aggregation operator under an identical protocol,
not a claim against the state of the art. Second, the result is not explained by an
absence of parallel paths: peptide graphs contain rings and branches, and a
typical molecule has hundreds of vertex pairs joined by $\ge 2$ walks of length
$\le 4$ (which the quotient ideal would collapse). Rather, the over-squashing is
\emph{mild}: these graphs are small (median $138$ nodes), so a few layers of
one-hop attention already reach far enough, and the multiplicity into any node
does not grow large enough to saturate the embedding. The multi-hop support that
is decisive under severe over-squashing yields little benefit when the squashing
is mild. Walk Attention's advantage is therefore \emph{regime-dependent}---large
when the squashing is severe (Tables~\ref{tab:main} and~\ref{tab:ring}), and
comparable to one-hop attention when it is mild---while on real data the operator
continues to train stably at moderate cost (about $1.6\times$ the wall-clock of
one-hop GAT, averaged over seeds, far below dense all-pairs attention).

\begin{table}[t]
\centering
\caption{LRGB Peptides-func, test Average Precision (mean $\pm$ std over 4
seeds; higher is better), under a shared small budget (four layers, width $80$,
no positional encodings) that isolates the aggregation operator; published
results under the standard tuned $500$k-parameter protocol are higher for every
architecture \cite{dwivedi2022lrgb,tonshoff2023where}. Walk Attention ties GAT
and exceeds GCN; on this mild-over-squashing benchmark the multi-hop support
yields no decisive gain (contrast Table~\ref{tab:main}).}
\label{tab:lrgb}
\begin{tabular}{lc}
\toprule
\textbf{Model} & \textbf{test AP} \\
\midrule
GCN                          & $0.492\pm0.011$ \\
GAT                          & $0.526\pm0.004$ \\
Walk Attention (ours)        & $0.526\pm0.007$ \\
\bottomrule
\end{tabular}
\end{table}

% ---------------------------------------------------------------------
\subsection{Cost}
% ---------------------------------------------------------------------

Walk Attention operates on the walk-reachability support, whose size is the
number of nonzero entries of $\Adj^{\,k}$. On the bottleneck family this is a
small fraction of $n^2$ ($1.5$--$2.6\%$ at the deepest range, depending on the
depth), so the layer is far
cheaper than a dense graph Transformer while retaining its long range. The
supports are precomputed once per graph topology and cached, so the path-algebra
computation does not enter the training loop.
